\documentclass[11.5pt,a4paper]{article}
\usepackage[utf8]{inputenc}
\usepackage[T1]{fontenc}
\usepackage[french]{babel}
\usepackage{geometry}
\geometry{top=2.5cm, bottom=2.5cm, left=2.5cm, right=2.5cm}

% --- Polices ---
\usepackage{helvet}
\renewcommand{\familydefault}{\sfdefault}

% --- Design & Couleurs ---
\usepackage[dvipsnames]{xcolor}
\usepackage{tcolorbox}
\tcbuselibrary{skins, breakable}
\definecolor{DeepBlue}{HTML}{154360}
\definecolor{AccentGold}{HTML}{D4AC0D}

\newtcolorbox{ConceptCle}[1]{
  colback=blue!5,
  colframe=DeepBlue,
  fonttitle=\bfseries,
  title=#1,
  arc=0mm,
  breakable
}

% --- Header/Footer ---
\usepackage{fancyhdr}
\pagestyle{fancy}
\fancyhf{}
\lhead{L'Affirmation de Soi : Pilier de Souffles Positifs}
\rhead{Expertise J.P. Perrin}
\rfoot{Page \thepage}

% --- Titre ---
\title{
    \Huge \textbf{L'Affirmation de Soi} \\
    \Large \textit{Du "Handicap du Non-Dit" au Leadership Authentique}
}
\author{Antigravity Intelligence System -- Rapport Consolidé}
\date{\today}

\begin{document}

\maketitle
\tableofcontents
\newpage

\section{Introduction : La Neurologie du Non-Dit}
L'incapacité à s'affirmer n'est pas un trait de caractère, c'est un mécanisme de défense neurologique. Le "non-dit" agit comme un poison lent qui sature les circuits de la récompense et active les zones de la douleur sociale.

\section{Diagnostic : Les Freins et l'Inhibition}
\subsection{Mécanismes Internes}
L'affirmation de soi est la capacité d'exprimer ses émotions tout en respectant l'autre. L'inhibition naît de trois peurs primaires :
\begin{itemize}
    \item \textbf{Peur du Jugement} : Activation du rejet social perçu comme une menace vitale.
    \item \textbf{Perfectionnisme} : L'impossibilité d'émettre une opinion "imparfaite".
    \item \textbf{Croyance Limitante} : "Dire non, c'est être agressif".
\end{itemize}

\subsection{Le Rôle Messager des Émotions}
\begin{ConceptCle}{L'Émotion comme Boussole}
Une émotion n'est pas un obstacle, c'est une donnée. La colère signale un besoin de justice, la tristesse un besoin de réconfort. Ne pas s'affirmer, c'est éteindre son propre tableau de bord.
\end{ConceptCle}

\section{De la Théorie à la Pratique : Tableaux de Transformation}
Basé sur les documents de référence (\textit{analyse\_finale\_affirmation\_leadership.md}), voici les horizons de progression.

\subsection{Matrice des Exercices et Bénéfices}
\begin{center}
\begin{tabular}{|p{4cm}|p{5cm}|p{5cm}|}
\hline
\textbf{Exercice} & \textbf{Objectif (Bénéfices)} & \textbf{Peines adressées} \\
\hline
Les 5 "Pourquoi" & Clarifier l'intention profonde. & Manque de sens initial. \\
\hline
Le "JE" assertif & Exprimer un besoin sans blâmer. & Agressivité réactive. \\
\hline
Défis Vision (Vidéo) & Travailler la posture et le regard. & Peur de l'exposition physique. \\
\hline
Cartographie & Définir son propre leadership. & Conformisme social. \\
\hline
\end{tabular}
\end{center}

\section{Le Leadership Authentique}
Le leadership n'est pas une question d'autorité, mais d'alignement intérieur. S'affirmer, c'est poser des limites saines qui permettent d'inspirer les autres par l'exemple, et non par la contrainte.

\section{Plan d'Action pour un Document Final "Pilier"}
Pour passer d'une ébauche à une version diffusable mondialement :
\begin{enumerate}
    \item \textbf{Standardisation scientifique} : Intégrer des sources neuroscientifiques post-2020.
    \item \textbf{Études de Cas} : Transformer les témoignages (ex: Paul le bénévole) en archétypes.
    \item \textbf{Mise en Forme} : Finalisation sous LaTeX Premium pour une impression professionnelle.
\end{enumerate}

\section{Conclusion}
L'affirmation de soi est le premier pas vers la liberté. En comprenant ses émotions comme des alliées, on transforme le handicap du silence en une force de leadership.

\end{document}
